We study rotor-fault tolerance for a simulated 5-inch FPV-class quadrotor (0.64 kg, thrust-to-weight 3.6) with sensor noise, a 10 ms motion-capture latency, parameter uncertainty and light turbulence. An online estimator of the four rotor effectiveness factors (filtered rigid-body regression, Kalman filter) feeds a bank of CUSUM tests for detection and isolation; a supervisor then either re-allocates thrust with the estimated effectiveness (partial loss) or, when a rotor can no longer carry its share of the weight, switches to a yaw-relaxed controller in the spirit of Mueller and D'Andrea (2014) that spins the vehicle about a tilted primary axis (gain-scheduled LQR on the reduced attitude). In 200 Monte-Carlo runs with 20-90 \% loss of effectiveness, all faults were detected (median delay 30 ms, 95th percentile 58 ms) and all were isolated to the correct rotor, with a median magnitude error of 1.5 percentage points; no false alarm occurred in 32 min of fault-free flight. With a 50 \% loss in hover the median peak deviation drops from 28 cm (passive) to 4.3 cm. After a complete rotor loss in hover the passive controller lost the vehicle in 32/32 runs (median 0.30 s), while the proposed pipeline recovered 32/32 with a median peak deviation of 38 cm and a spin of 23 rad/s; it also recovered 24/24 complete losses during a 3.3 m/s figure-eight. These are simulation results with assumed, not identified, parameters; in the robustness sweeps complete-loss recovery started to fail at an actuator delay of 24 ms (2 ms assumed), a motor time constant of 80 ms (25 ms assumed) and 3x the default turbulence. In 110 PX4 software-in-the-loop flights of a Gazebo x500, replayed offline with parameters read from the simulator's model files, the same diagnoser detected and isolated all 40 injected partial losses (19-48 \%, median delay 120 ms) and all 16 complete losses (median 132 ms), with no false alarm in 22 min of fault-free flight; PX4's own failure detector flagged none of the partial losses and the complete losses only after 664 ms, once the vehicle had flipped.
